% Part: counterfactuals
% Chapter: minimal-change-semantics
% Section: introduction

\documentclass[../../../include/open-logic-section]{subfiles}

\begin{document}

\olfileid{cnt}{min}{int}

\olsection{تمهيد}

قدّم ستالنيكر ولويس تحليلين للشرطيات المخالفة للواقع، كقولنا «لو حُك
عود الثقاب لاشتعل»، وأراد كل منهما بيان شروط الصدق الصحيحة لهذا الضرب
من الكلام. والأصل المشترك بينهما أن ننظر، عند تقويم الشرط، في العوالم
الممكنة التي لا تفارق العالم الفعلي إلا بأقل ما يكفي لصدق المقدّم؛ فإن
صدق التالي في تلك العوالم صدق الشرط المخالف للواقع. ولنفرض مثلًا أني
أمسك عودًا وعلبته، ولا أصنع في العالم الفعلي غير النظر إليهما وتقدير ما
يكون لو حككت العود. فأقل عالم تغييرًا أحك فيه العود هو الذي أختار فيه
الفعل ثم أحكه، ولا يتبدل معه أمر آخر: فلا أقفز في الهواء، ولا يشتعل
شعري بالحك، ولا تزول فجأة قوة أصابعي، ولا يفاجئني في اللحظة نفسها ماء
مسدس، وما أشبه ذلك. وفي هذا الإمكان البديل يشتعل العود؛ فيصدق القول إنه
لو حُك لاشتعل.

ويمكن أن يُكسى هذا التصور الحدسي دلالة صورية لمنطقيات الشرط المخالف
للواقع. رمز لويس إليه بـ«$\boxright$»، ورمز ستالنيكر بـ~«$>$»؛ ونحن
نستعمل~$\cif$ ونلحقه بمنطق القضايا رابطًا ثنائيًا. فيكون إلى جانب
!!{formula} التي صورتها~$!A \lif !B$ !!{formula} أخرى صورتها
$!A \cif !B$. وكالدلالة العلاقية لمنطق الموجهات، تقوم هذه الدلالة على
نماذج تُقوّم فيها !!{formula} عند العوالم. ويُحد شرط الإرضاء
$\mSat{\OLId{latin-looped}{M}}{!A \cif !B}[\OLId{latin-ordinary}{w}]$ بما يكون من~$\mSat{\OLId{latin-looped}{M}}{!A}[\OLId{latin-ordinary}{w}']$ ومن
$\mSat{\OLId{latin-looped}{M}}{!B}[\OLId{latin-ordinary}{w}']$ في بعض العوالم الأخرى~$\OLId{latin-ordinary}{w}'$. فأي~$\OLId{latin-ordinary}{w}'$ نختار؟ نختار،
بحسب التصور، أقرب عالم أو عوالم إلى~$\OLId{latin-ordinary}{w}$ يصدق فيها
$\mSat{\OLId{latin-looped}{M}}{!A}[\OLId{latin-ordinary}{w}']$؛ فلا بد إذن أن يحمل النموذج علاقة «قرب» أيضًا.

ومثّل لويس الحالات المخالفة للواقع تمثيلًا بيانيًا نافعًا: يجعل كل عالم
ممكن مركزًا لكرات متداخلة تضم عوالم أخرى، وتُرسم دوائر متحدة المركز.
فالعوالم الواقعة بين كرتين سواء في قربها من المركز؛ وما ضمته كرة داخلية
أقرب، وما لم تضمه إلا كرة تحيط بها أبعد.
\begin{center}
\begin{tikzpicture}[scale=.7]\small
  \spheresystem{5}
  \draw (0,0) node {$\OLId{latin-ordinary}{w}$};
  \propositionintersect{0}{4}{40}{3.5}
  {\tikzset{proposition/.append={smooth,tension=1.4}}
  \proposition[shift={(1.6,-1)}]{90}{1}{30}{4}}
  \path (1.5,3) node[below] {$\formula{B}$};
  \path (3,0) node {$\formula{A}$};
\end{tikzpicture}
\end{center}
وأقرب عوالم~$!A$ هي العوالم~$\OLId{latin-ordinary}{w}'$ التي تصدق فيها~$!A$ وتقع في أصغر
كرة تحيط بالعالم المركزي~$\OLId{latin-ordinary}{w}$، وهي المنطقة الرمادية. وعلى هذا التصور
تصدق~$!A \cif !B$ عند~$\OLId{latin-ordinary}{w}$ متى صدقت~$!B$ في كل عالم من أقرب
عوالم~$!A$.

\end{document}



